\section{Discussion}
\label{sec:discussion}

\subsection{Bounded verification evidence}
\label{subsec:discussion-bounded-verification}

RQ1 provides bounded evidence that no declared property was violated in the
completed valid-model configurations. The property suite covers distinct
obligations, including replay protection, destination authorization, decision
consistency, value preservation, quorum enforcement, terminal-state
irreversibility, and release after timeout. The completed base TLC
timeout-release checks are vacuous because those profiles disable timeout.
The timeout-enabled family supplies the relevant abort-branch evidence.

The evidential scope is determined by the explored bounds. Alloy completed all
three predefined profiles, whereas TLC-large was fully censored by the frozen
timeout policy. Those timeouts are not property violations; they delimit the
configurations for which bounded verification completed under the available
resource budget.

Accordingly, H1 supports only the absence of observed violations in completed
bounded configurations. It does not extend to unexplored states, censored
TLC-large executions, or an unbounded correctness claim.

TLC and Alloy provide complementary corroboration rather than interchangeable
proofs. Their different formal semantics and state representations prevent
agreement across completed profiles from implying identical explored spaces or
equivalent proofs.

\subsection{Mutation sensitivity and property adequacy}
\label{subsec:discussion-mutation}

RQ2 complements RQ1 by asking whether the declared properties react when
selected protocol controls are removed. All ten predefined scientific mutants
were detected by their designated target properties, producing a mutation score
of 1.0. This makes the valid-model results more informative because controlled
defects in replay protection, receipt-dependent credit, terminal decisions,
timeout release, and quorum enforcement trigger the expected properties.

Tool-specific encodings can express the same defect intention through different
properties. For example, commit after abort is associated with
\texttt{TerminalState\allowbreak Irreversibility} in Alloy and
\texttt{Decision\allowbreak Consistency} in TLA+. The mutant-level conclusion
therefore concerns targeted defect detection, not identical property
formulations across tools.

The perfect mutation score remains restricted to the predefined catalogue. It
does not establish completeness against arbitrary implementation defects,
modelling errors, or other cross-shard failure modes.

RQ1 and RQ2 thus support complementary claims: the declared properties hold in
the completed valid-model configurations, and the evaluated property suite is
sensitive to the selected protocol controls when they are deliberately
weakened.

\subsection{Implementation-model trace conformance}
\label{subsec:discussion-conformance}

RQ3 adds implementation-facing evidence by replaying observable Java traces
against the TLA+ abstraction. The valid catalogue was accepted and the
deliberately corrupted catalogue rejected as expected, connecting the formal
model with the implementation behaviors exercised by the frozen scenarios.

The negative controls strengthen this result because rejection was checked at
the expected diagnostic point rather than treated as a generic replay failure.
This provides evidence that deliberately introduced mismatches are localized at
the intended abstract-action boundary.

The result remains limited to the selected observable abstraction, scenario
catalogue, negative mutations, and frozen seeds. It therefore supports bounded
implementation-model trace conformance rather than a statement about every
possible Java execution.

In particular, this evidence must not be interpreted as a formal
refinement proof or as behavioral equivalence between the Java
implementation and the TLA+ specification. Establishing refinement
would require a stronger semantic relation covering the full relevant
transition systems, whereas the present experiment validates
correspondence for the predefined observable traces.

RQ3 therefore complements internal model checking and mutation analysis by
confronting the formal abstraction with observable implementation executions.

\subsection{Verification cost and scalability}
\label{subsec:discussion-cost}

RQ4 shows that bounded verification is constrained by computational cost as
well as logical outcome. Measurements increase within both formal tools, but
H4 remains unconfirmed in its preregistered nonlinear form.

For Alloy, all three ordered profiles completed, and median elapsed time and
memory increased monotonically from small to medium and large. The Spearman
association of 1.0 for both metrics is consistent with that ordering, but three
ordinal points cannot identify a functional growth law.

TLC exposes a sharper feasibility boundary. The small and medium profiles
completed with increases in elapsed time, memory, and explored state-space size,
whereas all seventy TLC-large measured runs reached the frozen 1800-second
timeout. Complete censoring precludes a completed-run median for the large
profile and a three-level characterization comparable to Alloy. The timeout is
therefore a valid censored observation, not an experimental failure, and
retaining it makes the evaluated verification boundary explicit.

Within the TLC medium-bound configurations, median time and memory varied
across normal, unanimity (three of three), replay, and
timeout-oriented cases, indicating that verification cost depends on behavioral
structure as well as nominal bound size. These differences do not identify a
universal cost ordering, but they show that enabled behavior is part of the
verification envelope.

The evidence supports only a descriptive conclusion that bounded verification
cost increases for the evaluated configurations. It establishes neither a
nonlinear nor an asymptotic growth law. TLC and Alloy also use different
semantics, exploration mechanisms, and internal representations, so absolute
time, memory, or state-space measurements cannot support intrinsic
tool-superiority claims.

\subsection{Cross-layer evidence triangulation}
\label{subsec:discussion-triangulation}

The four research questions are intended to provide complementary evidence
rather than four independent demonstrations of protocol correctness. Each
addresses a different weakness in the validation argument.

RQ1 evaluates whether the declared properties survive the bounded
interleavings explored for valid formal models. Because absence of a
counterexample does not demonstrate sensitivity to relevant defects, RQ2
deliberately weakens selected protocol controls and checks whether their
designated properties detect the resulting violations.

Formal verification and mutation detection still concern formal
representations. RQ3 adds implementation-facing evidence by checking bounded
conformance of observable Java traces against the TLA+ abstraction, together with
deliberately corrupted abstract-trace controls. This connects the formal layer with the
implementation without implying refinement or behavioral equivalence.

RQ4 makes the computational boundary explicit by characterizing verification
cost as the evaluated bounds become more demanding. The observed increase in
cost and complete TLC-large censoring show that the formal evidence is obtained
within a finite resource envelope.

Together, the layers form an evidence chain from bounded property verification
through mutation-based property validation and bounded implementation-model
trace conformance to verification-cost characterization. The chain is
complementary rather than cumulative proof: model checking remains bounded,
mutation analysis remains catalogue-dependent, trace replay remains limited to
the chosen abstraction and scenarios, and cost characterization remains tied to
the evaluated configurations. No layer removes the limitations of another.

The evidence chain consequently cannot be treated as additive assurance.
Bounded model checks, targeted formal mutants, replay of projected Java
scenarios, rejection of mutated abstract traces, and artifact reconstruction
answer different questions. The contribution of this case study is to
identify the resulting inference boundaries and concrete observational gaps,
not to establish a novel general-purpose verification workflow.

% Discusión negativa y límites de novedad; no reivindicar un nuevo algoritmo.
\subsection{When distinct evidence layers fail to support the same claim}
\label{subsec:observation-adequacy}

The case study yields a falsifiable interpretation of cross-layer evidence:
a successful check supports only the claim expressible by the state
representation, test input, and observation function actually used. This
matters even if all reported verifiers return their expected outcomes.

\paragraph{Ownership cannot be certified by amount-only observation.}
Consider a UTXO ledger state $L$ and a state $L'$ differing only in the
recipient of one existing output, with identical output amounts and shard
assignments. Let $V_s(L)$ be the sum of output amounts in shard $s$. Then
\[
 \forall s:\quad V_s(L)=V_s(L'),
 \qquad\text{although}\qquad L\ne L'.
\]
Any deterministic detector whose only input is the vector
$(V_s(L))_s$ returns the same result for both states. This is a basic
indistinguishability observation, not a new verification theorem. It explains
why ownership-preserving observations are necessary to assert recipient
correctness. The M02/M03 Java executions provide concrete examples of
wrong-recipient outcomes, while the G2 ledger-journal audit reports divergence on the same
beneficiary-distinguishing attempts. The empirical finding
is about what the original workload omitted, not a general weakness of G2.

\paragraph{Boundary correspondence is not internal refinement.}
V1 abstracts receipt consumption and destination credit into a common step,
while the Java operation emits intermediate debit/credit events and restores
state after some injected failures. V2.2 exposes additional intermediate
states, but the G2-P3 witness shows that its coordinator-decision transition
structure does not encode the examined Java debit-then-rollback history as a
step-preserving execution. No inference of internal refinement follows from
successful V1 boundary-trace replay. In contrast, the incompatibility of that
particular internal abstraction does not establish impossibility of a more
coarse-grained refinement relation.

\paragraph{Negative controls are not implementation failures.}
The 300 RQ3 positive traces are produced by ten deterministic Java-simulator
scenarios under thirty seeds. The 300 negative RQ3 cases are abstract-trace
mutations, and their rejection measures replay sensitivity to known
invalid abstract actions. Treating these controls as 300 Java defects would
inflate the implementation-facing evidence. Counting repeated seeds as
independent fault classes would likewise be pseudoreplication.

\paragraph{What the closest literature already establishes.}
MongoDB's distributed transaction study combines modular formal analysis
with implementation-facing model-derived tests
\cite{schultz2025mongodb}. Cirstea et al. validate Java-program traces
against TLA+; Mocket and TraceLink connect formal models and executable
systems with stronger automated test/trace support
\cite{cirstea2025traces,wang2023mocket,hackett2025tracelink}.
Model-guided Fuzzing investigates model-level coverage of distributed-system
implementations \cite{gulcan2025modelfuzz}. More specifically, formal
model-guided blockchain conformance testing already uses a deterministic
blockchain simulator \cite{drobnjakovic2025blockchains}. The revised 2026
OmniLink preprint addresses unmodified concurrent systems and trace-order
uncertainty \cite{hackett2026omnilink}.
Our experiments do not outperform or replicate these systems, and they cannot
support methodological priority for combining model checking, mutation, and
trace validation. Instead, they supply a precisely scoped \emph{negative
empirical case study}: missing workload variation, non-equivalent abstract
transitions, and artificial negative controls can coexist with apparently
complete verification results. A generalizable assessment framework would
require a separately evaluated detector or protocol implementation and
paired, causally interpretable comparisons; neither is claimed here.
